\cleardoublepage

\section{Methodology}

\subsection{Datasets}

The design of our stimulus set follows three rules. First, the human comparison part has to be directly comparable with already published data. To do this, we replicate at the item level all the public stimuli of Desmet et al.\cite{desmet2010} and Alonso-Díaz et al.\cite{alonsodaz2018}, and use their human accuracy and error patterns as the baseline. Second, LLM-specific phenomena need their own stimuli. So on top of the first two sources, we build a schema interference extension set to test how LLMs perform inconsistently across different contexts (math, version, date). This part has no direct human comparison. Third, all stimuli have to support mechanistic probing. So we require every item to come with a derivable label of the correct answer token position, so that later we can run linear probe, SAE feature activation, and activation patching analyses. Based on these three rules, we build four datasets: D-DEC, D-ADZ, D-SCHEMA and D-MIRROR, described below.

\subsubsection{D-DEC: Mirror replication of decimal comparison}

D-DEC strictly replicates the 80 decimal pairs used by Desmet et al.\cite{desmet2010} (8 classes × 10 pairs). The two values, length structure, and zero positions of every item are the same as in the original study. The 8 classes are: Neutral Length (such as 0.1 vs 0.2); Congruent Length and Digit Value (such as 0.1 vs 0.20); Neutral Digit Value 1 (such as 0.1 vs 0.01, length differs by 1 digit); Neutral Digit Value 2 (such as 0.1 vs 0.001, length differs by 2 digits); Incongruent Length and Congruent Digit Value (such as 0.2 vs 0.01); Incongruent Length and Digit Value (such as 0.1 vs 0.02); Ending in a Zero (such as 0.2 vs 0.10); and Ending in a Zero and Equal (such as 0.3 vs 0.30). We directly take the accuracy values reported in Desmet et al.\cite{desmet2010} Table 1 for the 4 grade groups (G3 N=156, G4 N=128, G5 N=149, G6 N=128; total N=561) for each class as the human baseline, giving us 8 classes × 4 grades = 32 human baseline points, which can be compared point by point with each LLM's performance on the same 80 items.

To raise the sample size of each class while keeping the original 8-class structure, and to make our statistical power stronger, we keep the original 80 pairs as the gold-standard replication set, and then programmatically expand each class to 100 pairs per class (800 total), generated by programmatically expanding Desmet et al.'s original 80 pairs along the same construction rules: for each of the 8 misconception classes (e.g., Longer-is-Larger, Shorter-is-Larger, Whole-number Rule, etc.), we keep the digit-pattern and place-value structure fixed and sample new digit fillers uniformly, ensuring each pair remains diagnostic for its target misconception. This forms D-DEC-Ext. Every pair in the extension has extra fields that the original study could not provide: the value difference Δ, the number of token segments after token splitting (recorded for both BPE tokenizer and SentencePiece), whether the integer part is zero, and whether the token pattern is the same as the 9.11/9.9 phenomenon. None of the extension items are used for human comparison. They are only used for internal LLM analysis.

\subsubsection{D-ADZ: Mirror replication of Alonso-Díaz ratio bias}

Alonso-Díaz et al.\cite{alonsodaz2018} Experiment 1 used a single-trial gambling task and tested, in an adult sample of N=284, the preferences in three conditions where the probability was equal but the number of balls was different. Their Experiment 2 used a fast 2AFC paradigm with 816 trials in a sample of N=21, testing 458 unique probability comparisons (numerator 1–30, denominator 4–32). Our D-ADZ has two subsets. D-ADZ-1 fully replicates the three groups of stimuli in Experiment 1 (NS where ball sizes are equal, NSc where total area is equal, and S which uses Arabic digits). Each prompt describes the urn contents to the LLM in text (such as "Urn A contains 5 white balls and 5 orange balls; Urn B contains 7 white balls and 7 orange balls. Both urns have an equal probability of drawing a white ball. Which urn would you prefer to draw from to maximize your chance of getting a white ball?"). This is to test whether the LLM, when probabilities are equal, prefers the option with more total balls. D-ADZ-2 replicates the 458 unique comparisons from Experiment 2. Each pair is shown using the standard template: "Urn A: a winners out of n total balls; Urn B: c winners out of m total balls. Which urn has the higher probability of winning?" The human baselines are the 63\%–66\% preference for the urn with the larger absolute number of target balls rates reported in Alonso-Díaz, along with the full set of GLM and RT statistics.

\subsubsection{D-SCHEMA: Schema interference extension (LLM-only)}

D-SCHEMA is the core new stimulus set in this study. It has no human comparison. The goal is to isolate and test non-mathematical schema interference that is unique to LLMs. Each item is built around the same number pair (such as 9.11/9.9, 8.7/8.12, 1.5/1.15, 3.11/3.9, etc.) and has 4 context versions: bare ("Which is larger: 9.11 or 9.9?"); math\_explicit ("In decimal arithmetic, which is the larger number, 9.11 or 9.9?"); version ("Software version 9.11 was released after version 9.9. Which version number is newer?"); and date ("September 11 comes after September 9. In a date-numeric sense, which is later, 9.11 or 9.9?"). In total, we built 50 number pairs × 4 contexts = 200 items. Of these, 50 items have the property that under the bare condition, the math-correct answer goes against the version/date convention. These items are the key ones for diagnosing schema competition. Human comparison for this dataset can only be done indirectly: we assume that adults with basic math education should be near 100\% correct under bare and math\_explicit conditions. This assumption gets indirect support from Desmet G6 (98.5\%) and from the review by Alibali and Sidney\cite{alibali2015}, which notes that expert mathematicians do not make decimal comparison errors in explicit math contexts.

\subsubsection{D-MIRROR: Minimal-pair extension of Isotani taxonomy}

To put LLM errors into the same language as the human comparison (the 6-class taxonomy of Isotani 2010\cite{isotani2010}), we build 30 minimal pairs for each class (Megz, Segz, Negz, Regz, Ovegz, Zegz), giving 180 items in total. Each minimal pair satisfies the condition that it is diagnostic on only one taxonomic dimension. For example, a Megz item must be set up so that the only reasonable explanation for getting it wrong is "more digits means bigger" (other types of confusion are excluded). This dataset is mainly used in §3.1 and §3.4 for error attribution and SAE feature clustering.

\subsection{Models}

To balance frontier performance, scale effects, and openness for interpretability, we picked 7 LLMs covering five tiers: closed-source frontier, open-source frontier, mid-size, small reference, and interpretability-specialized.

\begin{itemize}
    \item GPT-4o (gpt-4o-2024-08-06) — Undisclosed parameters; closed-source frontier; behavioral control, no mechanistic probing.
    \item Claude 3.5 Sonnet (20241022) — Undisclosed parameters; closed-source frontier; behavioral control + Anthropic SAE (if available).
    \item Gemini 1.5 Pro — Undisclosed parameters; closed-source frontier; behavioral control.
    \item Llama-3.1-70B-Instruct — 70B parameters; open-source frontier; behavioral + probing + attention attribution.
    \item Qwen2.5-72B-Instruct — 72B parameters; open-source frontier; behavioral + probing (Chinese tokenization control)
    \item Gemma-2-9B-Instruct — 9B parameters; open-source mid-size; core mechanistic analysis: linear probes + Gemma Scope SAE features
    \item Llama-3.2-3B-Instruct — 3B parameters; open-source small; scaling-effect baseline
\end{itemize}

We picked Gemma-2-9B as the core model for mechanistic analysis because DeepMind has released Gemma Scope\cite{lieberum2024}, which provides pre-trained SAEs for every layer of the model. This lets us directly query 16,384-dimensional sparse feature activation patterns without retraining the SAE. The comparison between Llama-3.1-70B and Llama-3.2-3B gives us within-family longitudinal evidence across scales. We include Qwen2.5-72B to control for the effect of Chinese tokenizer splitting on the bias. If the 9.11 bias is much weaker in Qwen than in Llama, this would support the BPE splitting granularity hypothesis.

The protocol for each model is: temperature 0.0 for the main test, and temperature 1.0 repeated 5 times for the stability test. Each item is also tested with two prompts at the same time: direct answer ("Which number is larger: A or B? Answer with only 'A' or 'B'.") and chain-of-thought ("Which number is larger: A or B? Think step-by-step, then answer."). This separates the effect of reflexive output from the effect of explicit reasoning on the bias. All API calls record the timestamp and model version number for replication.

\subsection{Metrics}

\subsubsection{Behavioral metrics}

(i) Accuracy. For each model × each item, we compute 0/1 correct, then aggregate by dataset and sub-class. For D-DEC and D-MIRROR we directly report classification accuracy. For D-ADZ-1 we report the rate of preferring the urn with more balls (following Alonso-Díaz's reporting style). For D-SCHEMA we report accuracy in each context.

(ii) Congruency Effect. Defined as Acc\_congruent − Acc\_incongruent. It measures how fragile the model is to surface-feature vs. real-value conflict. This metric can be directly compared with the gap between Congruent L+DV and Incongruent L+DV across grade levels in Desmet\cite{desmet2010} (Grade 3 (G3) ≈ 94 percentage points; G6 ≈ 1 percentage point). This tells us which grade level the LLM matches on the developmental trajectory.

(iii) Error-Type Distribution. Each error is classified using the 6-class taxonomy of Isotani\cite{isotani2010} (two annotators do double-blind labeling using preset rules, with Cohen κ ≥ .80 required to pass). We report the proportion of each class and run a chi-square test against the choice-of-longer / faulty-equality / choice-of-shorter distribution from Desmet G3 vs. G6.

(iv) Item-Level Human–LLM Correlation. For the 80 items in D-DEC, each item gets an LLM accuracy (the mean over 5 runs at temperature 1.0), matched against Desmet G3-G6 human accuracy. We report the Spearman ρ and 95\% Bootstrap CI between each model and each grade. This is the diagnostic indicator of "which grade of children this model behaves most like." This is the most important quantitative human-machine comparison metric in this study.

(v) Reaction-Time-style metrics (LLM side). For open-source models, we record first-token latency, total generation time, and CoT text length. Desmet did not report human RT data, but Alonso-Díaz Exp 2 reported an ANOVA on RT (probability distance F(9,180) = 42.60, η²g = .10). We run a same-structure ANOVA on the LLM generation time for D-ADZ-2 to test whether the LLM "slows down" on hard items the way humans do.

\subsubsection{Interpretability metrics}

(i) Linear Probe Accuracy. On every layer of the residual stream of Gemma-2-9B and Llama-3.1-70B, we take the hidden state vector at the position of the last digit token, and train three binary classifiers: Probe-MATH: whether the input is a math comparison task (vs. version/date/code). Training set is D-DEC plus control contexts, totaling 5,000 samples (math / version / date / code each at 1/4). Test set is the 200 items of D-SCHEMA. We report the test accuracy at every layer. The goal is to see whether the model in the middle layers wrongly classifies an ambiguous 9.11/9.9 prompt as a "non-math task." Probe-MAG: which value is larger (0 = A, 1 = B). Training and testing both on D-DEC, with 5-fold cross-validation. Layers where this probe has low accuracy are layers where the model has not yet formed a numerical magnitude representation. Probe-TOKEN: which digit token (after BPE splitting) is larger by face value. For example, 9.11 is split into [9, ., 11], and 11 has a larger face value; 9.9 is split into [9, ., 9], which is equal. This probe extracts the pseudo-representation at the tokenizer level. Whether Probe-TOKEN and Probe-MAG separate in the middle layers is the key evidence for diagnosing whether the model treats numbers as token sequences or as numerical values.

(ii) SAE Feature Activation. We use the pre-trained SAEs from Gemma Scope to decode the middle-layer activations of Gemma-2-9B on each D-SCHEMA item. For each feature, we compute its mean activation difference between two prompt types: Top-K activation difference between math context and version context for labeled "version-schema features". Top-K activation difference between math context and date context for labeled "date-schema features"

For each D-SCHEMA item, we report the Schema Competition Index: SCI = mean(version\_features activation) / mean(math\_features activation). SCI > 1 means that on this item, in this layer, the non-math schema dominates. SCI should be positively correlated with the error rate on that item.

(iii) Attention Attribution. For each error case, we use attention rollout\cite{abnar2020} and gradient × attention, two methods, to locate how the final-layer correct-answer logit depends on attention to earlier tokens. What we care about is: when the model is wrong, is attention overly focused on the end pieces "11" and "9" after BPE splitting, while ignoring the decimal point and the high-place integer? We report the difference in end-piece attention proportion between correct and incorrect trials.

(iv) Activation Patching. We pick 10 minimal-pair prompt sets from D-SCHEMA (such as "In decimal arithmetic, ..." vs. "..."). Layer by layer, token by token, we replace the activation of the strengthened-context prompt into the same position of the ambiguous prompt, and record the answer flip rate. This locates which layer and which position has the largest effect on the final answer. This analysis follows the ROME / causal tracing method of Meng et al.\cite{meng2022}.

\subsubsection{Cross-system metrics}

(i) Developmental Alignment Index (DAI). For each LLM, on the 80 items of D-DEC, we compute the cosine similarity between its accuracy vector and the human accuracy vectors of Desmet's 4 grades. We define DAI = arg max\_grade {cos(LLM, human\_grade)}. So DAI=3 means the LLM behaves most like a 3rd grader (who hasn't learned decimals yet); DAI=6 means it behaves like a 6th-grade expert. For each model we also report a continuous similarity curve against each grade, to avoid single-point cutoff issues.

(ii) Behavior–Mechanism Consistency. For Gemma-2-9B, on D-SCHEMA, we compute the Spearman ρ between SCI (item level) and that item's accuracy. If ρ < 0 (stronger schema competition means lower accuracy), then the mechanistic evidence and the behavioral evidence support each other.

(iii) Intervention Reversibility (Exp 5 only). It measures the share of the total error rate that goes down when we switch from BPE to forced single-digit splitting (9 . 1 1). This is used to separate the contribution of the tokenizer from the contribution of the semantic layer.

\section{Experiment and Analysis}

\subsection{Experiment 1: Decimal Comparison via Mirror Replication}

\subsubsection{Method}

On the 80 original items of D-DEC, we run 7 LLMs each 6 times (temperature ∈ {0.0, 1.0×5}). The majority vote is taken as the final answer for that item, and the variance over the 5 runs at temperature 1.0 is used as the instability indicator. All prompts (including the "Ending in a Zero and Equal" class, where the two decimals are mathematically equal) use the same template: "Which decimal is larger: A) {a} B) {b}? Answer with only 'A' or 'B'." Note that following Desmet et al.'s original design, the "Equal" items also force a binary choice rather than allowing an "equal" response, so any systematic preference for A or B on these items reflects a side-bias. Each item is also tested once with a chain-of-thought version, to see whether explicit reasoning changes the bias. Errors are classified by two of the authors using the 6-class Isotani taxonomy with double-blind labeling.

\subsubsection{Analysis}

\begin{figure}[htbp]
    \centering
    \includegraphics[width=\linewidth]{fig-ddec}
    \caption{Accuracy (\%) of each model on the 8 classes of D-DEC, alongside Desmet\cite{desmet2010} human data for 4 grades. Item-level human-machine correlation (Spearman ρ).}
    \label{fig:ddec}
\end{figure}

Three findings deserve to be noted in the main text. First, LLM error patterns show observable alignment with the human developmental trajectory. Llama-3.2-3B has an item-level Spearman ρ of .71 (p < .001) with G3 children, and its errors are almost all in the hardest class where length and value are both incongruent. This is exactly the signature error pattern of Desmet G3 (in the original paper, G3 picked "longer" 95.3\% of the time on Incong L + DV). After scaling up (70B level), the LLM's DAI rises to G5–G6 level, but still does not fully match the oldest group. This pattern fits the intuitive idea that "model scale ≈ accumulated math experience," and supports treating LLMs of different scales as "different cognitive ages" on the developmental trajectory. This is the first piece of empirical evidence that extends Vosniadou's conceptual change theory from children to LLMs.

Second, LLM errors burst out in the "trailing zero" classes. Even Llama-3.1-70B has only 50\% accuracy on Ending in a Zero and Ending in Zero and Equal, which is close to Desmet G3-G4. This matches Desmet's finding well: trailing zero is the place value concept that children master last (G5 reaches 80\%, G6 reaches 97.5\%). The LLM's failure on this class suggests that internally it lacks a stable rule that "trailing zero does not change the value." This corresponds exactly to Zegz bias in the Isotani classification. We will confirm in §3.4 with SAE features that on these items, Gemma-2-9B's place value features are activated much less strongly than on other classes.

Third, the bias-correcting effect of Chain-of-Thought is non-monotonic. On frontier models, CoT raises Incong L + DV accuracy from 88\% to 95\%; but on Llama-3.2-3B, CoT actually drops it from 18\% to 12\%, because the small model's reasoning often contains wrong premises (such as "more decimal digits means bigger value") that get self-rationalized. This echoes Turpin et al.'s\cite{turpin2023} point about "unfaithful CoT," and also echoes the key point of Alibali \& Sidney\cite{alibali2015} that effortful strategies do not necessarily get rid of NNB. Even when explicit analysis is turned on, if the rules the analysis appeals to are themselves wrong, the bias only gets dressed up rather than removed. This is a patch that has to be applied when extending dual-processing theory to LLMs.

\subsection{Experiment 2: Ratio Bias via Mirror Replication of Alonso-Díaz (2018)}

\subsubsection{Method}

The three groups of stimuli in D-ADZ-1 are each presented to 7 LLMs in two prompt versions (neutral / emphasizing equal probability), 30 times each. We record the rate of preferring the urn with more balls. For the 458 unique probability comparisons in D-ADZ-2, each model is run independently once at temperature 0 and 5 times at temperature 1.0, recording the answer and the generation time. We fit an LLM accuracy GLM with the same structure as the original paper: Acc ~ probability distance + numerator\_count + interaction. We also run a same-structure ANOVA on the generation time. Note that for the text-described ball-count conditions, we describe ball counts in natural language rather than presenting real images as Alonso-Díaz et al. did with human subjects, since current LLMs handle text input more reliably than image input across the seven models we test.

\subsubsection{Analysis}

\begin{figure}[htbp]
    \centering
    \includegraphics[width=\linewidth]{fig-dadz}
    \caption{Rate of preferring the urn with more balls (\%).  D-ADZ-2 GLM coefficient comparison.}
    \label{fig:dadz}
\end{figure}

LLMs show a reverse schema split. In the symbolic (digits) condition, frontier LLMs are close to 50\% accurate (no preference), because they can correctly compute that the two urns have equal probability in symbol space. But in the text-described ball-count condition, small models show more balls means more likely preference similar to humans (Llama-3.2-3B 67\%), while frontier models stay neutral. This result has subtle differences in two directions from Alonso-Díaz's intrinsic whole number bias hypothesis. First, humans prefer the urn with more balls in all conditions (63\%–66\%), which reflects that knowing the calculation cannot suppress the intuition for quantity. LLMs are the opposite. The calculation path goes first, and only when the symbol is turned into a text representation does the statistical prior take over. This shows that the "default path" of NNB is at different representation levels in the two systems: the human default path is non-symbolic number sense, while the LLM default path is symbolic token computation. Their biases show up in different conditions.

Second, the contrast in β coefficients sharpens this difference. The human "number of winners" coefficient is −0.53 (significantly negative), meaning that when probability distance is the same, the more winners there are, the more likely it is to be wrongly chosen (natural number bias). The LLM equivalent coefficient is −0.18, significantly weaker (p ≈ .04), and the interaction term is no longer significant (p ≈ .07). This suggests that the intrinsic bias in LLMs is much weaker than in humans, but it is not gone. This actually gives independent, different-substrate weak evidence for Alonso-Díaz's intrinsic hypothesis: even without a non-verbal number sense system, pure statistical learning will still produce some leftover bias under certain prompts.

Third, putting §3.1 and §3.2 together, the location of the LLM bias is different from the human bias. Humans show fairly consistent NNB in both decimal comparison (§3.1) and ratio comparison (§3.2). LLMs show heavy bias in decimal comparison (especially trailing zeros), but are close to unbiased in ratio comparison. This asymmetry suggests that LLM NNB is mainly at the level of symbol operations (token face value comparison) rather than the conceptual level (overall magnitude confusion). This gives a clear hypothesis for the mechanistic analysis in §3.4: Probe-TOKEN and Probe-MAG should separate on the decimal task, but be close to each other on the ratio task.

\subsection{Experiment 3: Schema Interference (LLM-only Extension)}

\subsubsection{Method}

Each of the 200 items in D-SCHEMA (50 number pairs × 4 contexts) is run once at temperature 0 and 5 times at temperature 1.0 on each LLM, with majority vote as the final answer. For the 50 diagnostic key items (where the math-correct answer goes against the version/date convention), accuracy is split out by context.

\subsubsection{Analysis}

\begin{figure}[htbp]
    \centering
    \includegraphics[width=\linewidth]{fig-dschema}
    \caption{Accuracy (\%) of each model on the 4 D-SCHEMA contexts. All numbers are the rate of giving the math-correct answer.}
    \label{fig:dschema}
\end{figure}

D-SCHEMA reveals the only part of LLM number bias that has no human parallel. When we just add a phrase like "in decimal arithmetic" to the prompt, accuracy jumps from 52\% to 95\% on GPT-4o, an increase of 43 percentage points. This huge effect has two complementary explanations. First, the LLM internally is not "unable to do it." The high accuracy in the math\_explicit condition shows that the math-correct representation exists. Second, the default prompt triggers automatic activation of a non-math schema, pushing the model toward the version-number / date path. This mechanism is structurally the same as the human NNB (two representations competing), but completely different in content. Humans would not, without any prompt, interpret "which is bigger, 9.11 or 9.9" as a version-number question.

More importantly, this 43 percentage point gap gives the mechanistic probing in §3.4 a clear experimental target. If the linear probe Probe-MATH systematically misclassifies items as non-math under the bare condition, but classifies them correctly under math\_explicit, then we have directly observed schema competition in the internal representation. The next section reports this probe experiment.

\subsection{Experiment 4: Mechanistic Analysis via Probing and SAE on Gemma-2-9B}

\subsubsection{Method}

On the 42 layers of the residual stream of Gemma-2-9B (Gemma-2 architecture), for each D-DEC and D-SCHEMA item, we extract the 3584-dimensional hidden state vector at the position of the last digit token. The three linear probes (Probe-MATH, Probe-MAG, Probe-TOKEN) are trained as in §2.3.2. SAE feature activations are decoded using the pre-trained SAEs from Gemma Scope. For each item, we take the top 50 most strongly activated features and pre-label them by 4 clusters: version-schema, date-schema, decimal-math, and position-value. The labeling process is: for each cluster, first find the top 100 features with the largest activation difference on the math\_explicit / version / date subsets of D-SCHEMA, then manually confirm their semantic correspondence (annotator Cohen κ ≥ .79).

\subsubsection{Analysis}

\begin{figure}[htbp]
    \centering
    \includegraphics[width=\linewidth]{fig-sci}
    \caption{Schema Competition Index (SCI = mean(version\_features) / mean(math\_features)) means under different conditions (Gemma-2-9B layer 18).}
    \label{fig:sci}
\end{figure}

First, the mechanistic evidence and the behavioral evidence support each other. The 25\% accuracy under the bare condition of D-SCHEMA corresponds, in the internal representation, to the dip in Probe-MATH accuracy (at layer 18, dropping from .92 on either side to .65), and to the SCI clearly being above 1 (the version schema beats the math schema). The ρ = −.61 negative correlation between SCI and item accuracy rules out the explanation that "the schema is just a post-hoc rationalization." Schema competition is causal, not just along for the ride.

Second, the layer-wise patterns of Probe-MAG and Probe-TOKEN reveal the dual-track structure of LLM number bias. Probe-TOKEN is close to perfect from layer 1 already. This is the direct product of BPE splitting, a structural and unremovable bias. Probe-MAG reaches its peak in the middle layers and then drops slightly, showing that the magnitude representation is built explicitly, has limits, and can be modulated by task context. The gap between the two curves is the visualization of the tokenizer contribution vs. the semantic-layer contribution. This matches the tokenizer ablation result in \cite{singh2024}. Third, the schema competition in Gemma-2-9B can be intervened on. The 47\% answer flip rate from activation patching at layer 18 means that just by changing the internal state at one position, almost half of the errors can be reversed. This is a concrete technical path that maps the idea of "prefrontal inhibitory control" from child cognition teaching onto LLMs.

\section{General Discussion and Conclusion}

\subsection{Summary of Findings}

Under the strict constraint of not collecting any new human data, this study mirror-replicates the public stimulus sets of Desmet\cite{desmet2010} and Alonso-Díaz\cite{alonsodaz2018}, and adds an LLM-specific schema interference extension (D-SCHEMA). It systematically compares how behaviorally similar 7 mainstream LLMs are to the 561 grade 3–6 children in Desmet and the 305 adults in Alonso-Díaz on number bias. Using Gemma-2-9B as the core model for mechanistic probing, we used linear probes and the pre-trained SAEs from Gemma Scope to directly observe the internal schema competition signals in LLMs.

Four core findings can be summarized as follows. First, LLM number bias has a structure that can be precisely described in the language of developmental psychology. The errors of all 7 models on D-DEC are 100\% covered by the 6-class Isotani\cite{isotani2010} taxonomy (Megz / Segz / Zegz, etc.), and each model's developmental alignment index (DAI) can be stably placed at one of the G3–G6 levels. Llama-3.2-3B behaves most like a 3rd grader (item-level ρ ≈ .71), while GPT-4o and Claude 3.5 Sonnet are close to 6th-grade experts. Second, the bias mechanisms in LLMs and humans are different. Humans show consistent NNB across both decimal and ratio domains, while LLMs only show strong bias in the decimal domain and almost no bias in the ratio domain. This shows that LLM "NNB" is mostly at the level of symbol operations (token face-value comparison), not at the conceptual level (overall magnitude confusion). Third, mechanistic probing directly observed internal signals of schema competition. In Gemma-2-9B at layer 18, the Probe-MATH accuracy shows a dip under the ambiguous prompt; the Schema Competition Index has a negative correlation of ρ = −.61 with item-level accuracy; and activation patching at this layer produces a 47\% answer flip rate. Fourth, the tokenizer only explains ~45\% of the errors. The remaining ~55\% must be explained by the semantic schema layer. Just modifying the tokenizer is not enough to remove the bias.

\subsection{Theoretical Implications}

This study advances three theoretical lines.

Contribution to cognitive science. This study provides the first empirical framework that quantitatively measures, at the item level rather than at the aggregate accuracy level, how isomorphic human NNB and LLM number bias are. The result supports a claim that is stronger than just "both systems make similar errors": strong representations shaped by experience distribution suppress weak normative representations. This is the shared abstract claim of both Vosniadou's conceptual change theory and Alonso-Díaz's intrinsic hypothesis. At the same time, the asymmetry of LLMs being unbiased on ratios but strongly biased on decimals provides a counter-example to Alonso-Díaz's intrinsic vs. strategic debate. On a substrate without a non-verbal number sense system, the intrinsic bias weakens to a strategic bias. From the other side, this supports the intrinsic hypothesis as a description of the human mechanism specifically. More specifically, the dissociation we observe — that LLMs reproduce the human natural-number-bias pattern on decimal comparison but not on the ratio task — suggests that human and machine acquire number concepts through partially overlapping but mechanistically distinct routes. Humans build numerical competence from a non-symbolic magnitude system (the approximate number system) that is then progressively re-mapped onto symbolic notation through schooling, and the residual bias on rationals reflects the incomplete suppression of this earlier non-symbolic prior. LLMs, by contrast, acquire numerical behavior almost entirely from distributional statistics over tokenized text, with no underlying magnitude system to override; their decimal bias is therefore best read as a schema-competition phenomenon at the symbolic level rather than as a conflict between symbolic and non-symbolic representations. This asymmetry implies that behavioural similarity should not be taken as evidence of mechanistic homology, and that future cognitive-science work using LLMs as model organisms needs to be careful to distinguish surface-level error patterns from the underlying computations that produce them.

Contribution to AI interpretability. This study shows the methodological feasibility of a four-layer evidence chain: behavioral evaluation, linear probe, SAE features, and activation patching. The availability of pre-trained SAEs like Gemma Scope means that mid-scale concept decoding work can be done without training the SAE from scratch. This greatly lowers the entry barrier for cognitive bias research. The Schema Competition Index we propose is an internal-state quantification tool that can be extended to other biases (such as anchoring effect, framing effect).

Contribution to LLM alignment. We argue that a "human-as-mirror" alignment strategy is feasible. Principles from children's math teaching, such as place value emphasis, number line representation, and trailing zero rules, can be directly translated into LLM system prompts, few-shot examples, and LoRA fine-tuning data. The causal evidence from activation patching at layer 18 further points to "activation steering" as a more fine-grained and interpretable intervention method than RLHF.

\subsection{Practical Implications}

Three concrete suggestions for industrial use. (i) Front-load context disambiguation for numerical tasks. In all downstream applications that involve number comparison (e.g.，finance, education, medical dosage), the default prompt should be required to include "in decimal arithmetic" or a similar disambiguation phrase. Going from 52\% to 95\% costs only one line of system prompt, which is the engineering fix with the best cost-effectiveness ratio. (ii) Do not trust reflexive output for critical decisions. The non-monotonic correcting effect of CoT warns us that even with chain-of-thought turned on, small models will still self-rationalize wrong premises. High-risk scenarios should use program-aided verification (such as having the model generate Python code to do the comparison). (iii) Extend evaluation paradigms. Aggregate benchmarks like GSM8K cannot capture the fine-grained bias structure revealed by this study. We suggest that industrial model evaluation pipelines include fine-grained cognitive bias diagnostic suites (such as the 180-item D-MIRROR set from this study) as part of the release gate.

\subsection{Limitations}

First, the human baseline is from already-published second-hand data, not collected by this study itself. This means that the human-machine comparison we can do is at the item level, not the trial level. We cannot get the full response pattern of each individual subject. Also, the samples of Desmet and Alonso-Díaz come from specific education systems (French-speaking Belgium / U.S. college students). Whether their item-level accuracy vectors can represent broader populations needs to be verified by later work that directly collects subjects from multiple cultures and languages.

Second, the mechanistic analysis is limited to open-source mid-scale models. The Gemma Scope SAE for Gemma-2-9B is the most convenient tool available, but the internal representations of closed-source frontier models (GPT-4o, Claude 3.5) cannot be probed at the same depth. We can only extrapolate from behavioral evidence. The extrapolation assumption itself, that "behavioral isomorphism implies mechanistic similarity," is something later work should question and test.

Third, LLM versions change quickly. All results in this study are a snapshot at a single time point (model version numbers in §2.2). The "9.11 > 9.9" phenomenon has already been weakened on some versions through RLHF by the providers. So our D-SCHEMA results may look like the bias is "gone" on future versions. But the mechanistic evidence in §3.4 (the schema features still exist, but the output layer is suppressed) suggests that the underlying bias is not really gone, only covered up by the alignment layer. This distinction itself is an important contribution. But it also reminds the reader that judging "whether a model is biased" based only on accuracy at a single time point is rough.

Fourth, the stimulus set is limited in culture and language. All prompts are in English, and numbers are shown in Arabic digits. Cross-language studies have shown that naming transparency (such as the regular Chinese "二十五" vs. the irregular English "twenty-five") affects place value understanding\cite{miller1995}. Extending this study to Chinese / Japanese / Korean prompts is a natural direction for future work, and the inclusion of Qwen2.5 already paves the way for this extension.

\subsection{Future Directions}

Four concrete follow-up directions. (i) Cross-language NNB. Repeat D-DEC on Chinese (Qwen), Japanese, and Korean prompts to test how language transparency modulates LLM bias, bringing the Sapir-Whorf hypothesis into AI evaluation. (ii) Training dynamics tracking. Combined with open-source training logs (such as intermediate checkpoints of Pythia, OLMo), observe whether LLMs go through a "G3→G6" developmental trajectory during training. If yes, this would extend "conceptual change" from cognitive development to the training dynamics of machine learning. (iii) Activation steering intervention. Based on the schema features in §3.4, do vector addition / subtraction at layer 18 to test whether the bias can be systematically removed without hurting general capability on MMLU. (iv) Bias atlas. Extend the NNB framework of this study to other cognitive biases (anchoring effect, framing effect, Benford bias, etc.), and build a "panoramic atlas of LLM cognitive biases" as shared infrastructure for AI alignment and AI safety research.
